% Replaces review4_theory_proofs.tex; Appendix A.1 uses generic ell,w.
\subsection{Sharp reversal for nonnegative ratio risks}
\paragraph{Proof of Proposition~\ref{prop:lowdemand}.}
Since $\ell=\rho w$ on $U$, its error-to-exposure ratio is $\rho$
and its excess mass is $(\rho-r)M_U$. The combined excess is
$-B+(\rho-r)M_U$, proving the first assertion and the coverage increments.

For sharpness, take $k,p,q>0$ with $k+p+q=1$, $k>c_0$, $q>p$,
and $0<\varepsilon<p/q$. Set the conditional moments to
\[
\begin{array}{c|ccc}
 &K&U&V\\\hline
 \Pr(X)&k&q&p\\
 w&\dfrac{q(\rho-r)\varepsilon}{rk}&\varepsilon&1\\
 \ell&0&\rho\varepsilon&r+\dfrac qp(\rho-r)\varepsilon
\end{array}
\]
They can be realized simply by deterministic $(L,W)=(\ell,w)$.
Their excesses are
\[
 m_K=-\frac{q(\rho-r)\varepsilon}{k},\quad
 m_U=(\rho-r)\varepsilon,\quad
 m_V=\frac qp(\rho-r)\varepsilon.
\]
Every optimum accepts $K$, since doing so increases either utility
and relaxes the constraint. Its budget is $B=q(\rho-r)\varepsilon$.
Once $K$ is accepted, feasibility reduces to $a_U+a_V\leq1$.
Row utility is maximized uniquely at $(a_U,a_V)=(1,0)$ since $q>p$;
exposure utility is maximized uniquely at $(0,1)$ since $q\varepsilon<p$.
These arguments optimize over all fractional acceptance rules.
Both policies have ratio risk $r$ and row coverage above $c_0$.
With $T=p+q\rho\varepsilon/r$,
\[
 \Delta c=q-p,\qquad
 d_W(a_R)=\frac{q\rho\varepsilon}{rp+q\rho\varepsilon},\qquad
 \Delta d_W=\frac{r(p-q\varepsilon)}{rp+q\rho\varepsilon}.
\]
Take $k>c_0$ sufficiently near $c_0$ and then $p>0$ sufficiently
small that $q-p=1-k-2p>1-c_0-\eta$; for these masses,
$\varepsilon\downarrow0$ gives $\Delta d_W\to1$.
The opposite inequalities follow from $c(a_R)\leq1$,
$c(a_W)\geq c_0$, and $0\leq d_W\leq1$.

For $0<r<1$, choose $\rho=1$ and scale all moments by
$h=1/\max(1,w_K)$, which leaves risks, both coverages, and both
optimizers unchanged. Then $0\leq\ell\leq w\leq1$.
At each context, put probabilities $\ell,w-\ell,1-w$ on
$(L,W)=(1,1),(0,1),(0,0)$. This realizes the same conditional moments
with nested binary events, proving bounded-loss sharpness.
For the deterministic WAPE specialization before scaling, take
$Y=(w_K,\varepsilon,1)$ and
$f=(w_K,0,(1-r)(1-q\varepsilon/p))$. These nonnegative forecasts
realize precisely the $\rho=1$ moments. \hfill$\square$

\paragraph{Scope of the realizability condition.}
The displayed moment family is a sufficient condition for a fixed
loss/exposure model class to inherit the sharp bounds. In this
three-context family the strict inequalities $q>p$ and
$q\varepsilon<p$ are also necessary for these two distinct, unique
optima. They are not necessary conditions for reversal in arbitrary
larger classes. Ratio form, nonnegativity, and integrability alone
do not make every specified metric class rich enough: constant
$w$ makes exposure and row utility proportional. For a precision-like
error ratio, $L=\ind\{\widehat Y=1,Y=0\}$ and
$W=\ind\{\widehat Y=1\}$; for a recall-like miss ratio,
$L=\ind\{Y=1,\widehat Y=0\}$ and $W=\ind\{Y=1\}$.
Both have $L\leq W$, but realization requires acceptance to be
chosen using information for which the joint event probabilities
can take the three specified values. A deterministic prediction
included in $X$ restricts these moments. Recall here is conditional
on accepted cases; a metric counting rejected positives in its
fixed denominator is a different constraint.

\paragraph{A margin-dependent bound.}
For alternatives $K\cup U$ and $K\cup V$, suppose $\ell=\rho w$
on $U$ and $U$ exhausts the common slack:
$B=(\rho-r)M_U>0$. If $K\cup V$ satisfies the cap, $M_V>0$,
and $R(\ind_V)-r\geq\delta>0$, then
\[
 d_W(\ind_{K\cup V})-d_W(\ind_{K\cup U})
 \leq\frac{M_U}{T}\left(\frac{\rho-r}{\delta}-1\right).
\]
Indeed, $\delta M_V\leq\E[\ind_Vm_r]\leq B$.
If both alternatives exhaust the same budget, the exact identity is
\[
 \Delta d_W=\frac{M_U}{T}
 \frac{\rho-R(\ind_V)}{R(\ind_V)-r}.
\]
The construction has $M_U/T\to0$ but $R(\ind_V)-r\to0$ as well;
this is why exposure share alone gives no vanishing bound.
For WAPE near-zero forecasts satisfy
$|\E[\ind_U e]-M_U|\leq\E[\ind_Uf]$, and hence
$|R(\ind_U)-1|\leq\E[\ind_Uf]/M_U$.

\paragraph{Finite perturbations of estimated rankings.}
Let $\widehat\ell=\ell+\delta_\ell$,
$\widehat w=w+\delta_w$, where $w>0$ and
$|\delta_w|\leq\beta w$ for some $0\leq\beta<1$. Direct subtraction gives
\[
 \left|\frac{\widehat\ell}{\widehat w}-\frac\ell w\right|
 =\frac{|w\delta_\ell-\ell\delta_w|}{w|w+\delta_w|}
 \leq\frac{|\delta_\ell|+(\ell/w)|\delta_w|}{(1-\beta)w}.
\]
In contrast, $|\widehat m_r-m_r|\leq|\delta_\ell|+r|\delta_w|$.
The local derivatives with respect to $w$ are $-\ell/w^2$ and $-r$;
the ratio's sensitivity with respect to $\ell$ is $1/w$.
These are sensitivity statements, not a proof that denominator
noise alone causes the observed empirical ordering.
For fixed $T>0$, write $b=(1-\lambda)/T$ and
$s_\lambda=(\ell-rw)/(\lambda+bw)$. When $0<\lambda\leq1$ and
$w,\widehat w\geq0$, the denominator is at least $\lambda$, and
\[
 \left|\widehat s_\lambda-s_\lambda\right|
 \leq\frac{|\delta_\ell|+r|\delta_w|}{\lambda}
 +\frac{|m_r|b|\delta_w|}{\lambda^2}.
\]
Its derivatives satisfy
\[
 \frac{\partial s_\lambda}{\partial w}
 =-\frac{r\lambda+b\ell}{(\lambda+bw)^2},\qquad
 \frac{\partial s_\lambda}{\partial\ell}
 =\frac{1}{\lambda+bw}\leq\frac1\lambda.
\]
Thus positive row utility removes the small-denominator singularity
for fixed bounded moments. It does not remove fitting error,
calibration error, or distribution shift. A floor $\max(\widehat w,c)$
likewise bounds denominator sensitivity, while changing the
unregularized population ranking.
